\chapter{Több-adatkészletes architektúra}

\section{Miért vált szükségessé a dataset-regiszter?}

Az egyetlen adatbázisos modell korai prototípushoz még elegendő volt, de amint a projekt egyszerre kezdett demo-, smoke-test- és kutatási célokat szolgálni, világossá vált, hogy a különböző minták szétválasztása nélkül a rendszer nehezen reprodukálható. Ugyanaz a frontend, ugyanaz a Rust futtatókörnyezet és ugyanaz a Python pipeline más eredményt adhat, ha eltérő bemeneti korpuszon fut. Ezt a különbséget nem rejtett konfigurációs fájlokban, hanem explicit dataset-regiszterben érdemes kezelni.

A dataset-regiszter bevezetésének szükségességét konkrét igényváltozás indokolta:

\begin{itemize}
    \item \textbf{Smoke-teszt igény}: a fejlesztői pipeline-t kis adathalmazzal kell tudni validálni, anélkül hogy az éles adatbázist módosítani kellene;
    \item \textbf{Kutatási szétválasztás}: a Flex Queue és a Solo/Duo Queue strukturálisan eltérő gráfot hoz létre; az eredmények összekeveredése módszertanilag hibás lenne;
    \item \textbf{Reprodukálhatósági garancia}: ha a fő kutatási dataset változik, a régi dataset-állapot is ellenőrizhető maradjon;
    \item \textbf{Pipeline-átjárhatóság}: a Rust analitika, a Python klaszterpipeline és a Node/Express API mind ugyanarra a dataset-re álljon rá, ne legyen ``backend igazság'' és ``Rust igazság'' egymás mellett.
\end{itemize}

Az \textbf{elvetett alternatíva} a konfigurációs fájlon alapuló manuális váltás lett volna: a felhasználó kézzel módosít egy útvonalat, és a rendszer azt használja. Ez nem skálázódik, mert a több állapot egyszerre való megtartása nem lehetséges, és az összekeveredési hibák nehezen detektálhatók.

A \texttt{backend/data/datasets.json} pontosan ezt a feladatot látja el. A dataset fogalma a jelenlegi rendszerben nem csak annyit jelent, hogy ``másik mappa'', hanem egy teljes, együtt mozgó adategységet:

\begin{itemize}
    \item nyers adatbázissal,
    \item finomított adatbázissal,
    \item meccs-JSON könyvtárral,
    \item cache-könyvtárral,
    \item megjeleníthető metaadatokkal.
\end{itemize}

\section{A jelenlegi regiszterben szereplő adatkészletek}

A dolgozat írásakor ellenőrzött \texttt{backend/data/datasets.json} regiszterben négy dataset szerepelt:

\begin{itemize}
    \item \texttt{default} - az eredeti telepítésből migrált, vegyes queue-tartalmú alapkorpusz;
    \item \texttt{smoke-test-dataset} - gyors infrastruktúratesztelésre szánt átmeneti adatkészlet;
    \item \texttt{flexset} - EUNE Master$+$ Flex Queue meccsek, kifejezetten premade és ismétlődő csapattársi kapcsolatok tanulmányozására;
    \item \texttt{soloq} - EUNE Master$+$ Solo/Duo Queue meccsek, kontroll- és teljesítmény-alapvonal-gyűjtés céljából.
\end{itemize}

Az \textbf{aktív dataset a \texttt{soloq}}. Ez a döntés az adatbázis-regiszterben is rögzített: a \texttt{currentDatasetId} mező értéke \texttt{"soloq"}.

A négy dataset jelenléte nem véletlen felhalmozódás, hanem tudatos kutatásstratégia eredménye. A \texttt{default} megőrzi az eredeti, heterogén korpuszt és a korábbi fejlesztési előtörténetet. A \texttt{smoke-test-dataset} a pipeline gyors ellenőrzésére szolgál anélkül, hogy a főadatokat érintené. A \texttt{flexset} és a \texttt{soloq} egymással párban értelmezhetők: az előbbi a premade és csoportos játékhoz kötődő társas jeleket erősíti, az utóbbi pedig ugyanannak a szervernek és rangsornak egyénibb, kevésbé premade-fókuszú meccstörténetét képviseli.

\begin{table}[H]
\centering
\caption{A négy regisztrált dataset összefoglalója}
\begin{tabularx}{\textwidth}{>{\raggedright\arraybackslash}p{3.0cm} >{\raggedright\arraybackslash}p{3.0cm} X}
\toprule
\textbf{Dataset ID} & \textbf{Szerepkör} & \textbf{Jellemzők} \\
\midrule
\texttt{default} & Alap/migrált & Vegyes queue, heterogén tartalom \\
\texttt{smoke-test} & Pipeline-teszt & Átmeneti, infrastruktúra-validálásra \\
\texttt{flexset} & Fő kutatási dataset & EUNE Master$+$ Flex, premade-fókusz \\
\texttt{soloq} & Kontroll/baseline & EUNE Master$+$ Solo/Duo, teljesítmény-alap \\
\bottomrule
\end{tabularx}
\end{table}

A \texttt{flexset} esetében a regiszter metaadata tartalmazza a gyűjtési stratégia dokumentumát és a collector konfigurációs fájl hivatkozását (\texttt{apex\allowbreak-flex\allowbreak-collector.json}). A \texttt{soloq} hasonlóan strukturált: külön collector konfigurációs fájllal és gyűjtési stratégia-dokumentummal rendelkezik. Ez a párhuzamos konfiguráció kifejezetten mutatja, hogy a két dataset nem eseti gyűjtés, hanem tudatosan eltérő stratégiával felépített, összehasonlítható korpusz.

\begin{table}[H]
\centering
\caption{A dataset-regiszter funkcionális szerepei}
\begin{tabularx}{\textwidth}{>{\raggedright\arraybackslash}p{4.2cm}X}
\toprule
\textbf{Funkció} & \textbf{Jelentőség} \\
\midrule
Aktív dataset kijelölése & Meghatározza, hogy a backend és a frontend mely adatkorpuszt használja alapértelmezetten \\
Eltérő tárolási útvonalak & Megakadályozza, hogy a teszt- és főadatok összekeveredjenek \\
Külön cache-ek & A bird's-eye és egyéb előre generált artifaktumok elkülönülnek \\
Kísérleti összevethetőség & Azonos pipeline különböző mintákon is futtatható \\
\bottomrule
\end{tabularx}
\end{table}

\section{A dataset-választás propagálása a backend és a Rust felé}

A dataset-ek architektúrájának egyik legfontosabb technikai előnye az, hogy a választás nem áll meg a Node-szinten. A backend a kiválasztott datasetből környezeti változókat és útvonalakat képez a Rust futtatókörnyezet felé is, például:

\begin{itemize}
    \item \texttt{DATASET\_ID},
    \item \texttt{GRAPH\_DB\_PATH},
    \item \texttt{PATHFINDER\_MATCH\_DIR},
    \item \texttt{PATHFINDER\_BIRDSEYE\_CACHE\_DIR}.
\end{itemize}

Ez építészetileg nagyon fontos döntés. Nincs külön ``frontend igazság'', ``backend igazság'' és ``Rust igazság'' ugyanarról az adatállapotról: a dataset-választás a teljes végrehajtási láncon végigmegy.

A \texttt{server.js} szintjén ez úgy valósul meg, hogy a szerver a kiválasztott datasetből egységes környezeti változókészletet képez, például \texttt{DATASET\_ID}, \texttt{GRAPH\_DB\_PATH}, \texttt{PATHFINDER\_MATCH\_DIR} és \texttt{PATHFINDER\_BIRDSEYE\_CACHE\_DIR} formában. A Rust runtime ezeket fogyasztja, így a teljes végrehajtási lánc ugyanarra az adatvilágra áll rá. Ez lényeges a védhetőség szempontjából: a frontend által mutatott eredmény nem különálló demonstrációs adatból, hanem ugyanabból a datasetből jön, amelyen a kutatási analitika is fut.
